\section{Taller resuelto: Consumo efectivo de motores distintos}

\noindent\textbf{Nivel:} Intermedio.

\apartado{1. Enunciado}{Dos motores entregan 24 y 19 kN con consumos específicos 16 y 19 g/(kN s). Calcule gasto total, cT efectivo y combustible en 32 minutos. Compare con promediar sin ponderar.}

\apartado{2. Datos}{T1=24000 N, T2=19000 N; consumos másicos referidos a empuje.}

\apartado{3. Incógnitas}{Las magnitudes y comprobaciones solicitadas en el caso.}

\apartado{4. Sistema de referencia}{Empujes paralelos; suma escalar longitudinal.}

\apartado{5. Hipótesis físicas}{Ambos puntos de funcionamiento permanecen constantes durante el intervalo.}

\apartado{6. Modelo matemático}{Consumo efectivo de motores distintos}

\apartado{7. Ecuaciones simbólicas}{Se parte de los balances y definiciones siguientes.}

\begin{equation}\dot m_f=c_{T1}T_1+c_{T2}T_2\end{equation}

\begin{equation}c_{T,ef}=\dot m_f/(T_1+T_2)\end{equation}

\begin{equation}\Delta m=\dot m_ft\end{equation}

\apartado{8. Desarrollo algebraico}{Las operaciones se realizan en este orden, conservando las condiciones del modelo.}

\begin{enumerate}

\item Se calcula el gasto de cada motor con su propio consumo.

\item Se suman gastos y empujes y se obtiene la razón efectiva.

\item La media aritmética sólo sería equivalente si los empujes fuesen iguales.

\end{enumerate}

\apartado{9. Sustitución numérica}{Se conservan los resultados intermedios sin redondear y se muestran todas las operaciones escalares disponibles.}

\begin{enumerate}

\item Gasto total.

\[\fitmath{y_{1}=1.6\times10^{-5}\cdot 24000+1.9\times10^{-5}\cdot 19000\approx 0.745\,\mathrm{kg/s}}\]

\item Consumo efectivo.

\[\fitmath{y_{2}=\frac{1.6\times10^{-5}\cdot 24000+1.9\times10^{-5}\cdot 19000}{24000+19000}\approx 1.732558\times10^{-5}\,\mathrm{kg/(N\,s)}}\]

\item Combustible.

\[\fitmath{y_{3}=\left(1.6\times10^{-5}\cdot 24000+1.9\times10^{-5}\cdot 19000\right)\cdot 32\cdot 60\approx 1430.4\,\mathrm{kg}}\]

\end{enumerate}


\apartado{10. Resultados con unidades}{Los resultados pertenecen al modelo y al intervalo de datos declarados.}

\begin{center}\begin{tabular}{lr}\toprule Magnitud&Resultado\\\midrule

Gasto total & $0.745\,\mathrm{kg\,s^{-1}}$\\

Consumo efectivo & $1.732558\times10^{-5}\,\mathrm{kg\,N^{-1}\,s^{-1}}$\\

Combustible & $1430.4\,\mathrm{kg}$\\

\bottomrule\end{tabular}\end{center}

\apartado{11. Interpretación}{El consumo efectivo es una media ponderada por empuje.}

\apartado{12. Comprobación independiente y alcance}{Calcular en gramos por segundo y comprobar que multiplicar cT efectivo por empuje total recupera el mismo gasto.}

\apartado{13. Limitaciones}{No incluye asimetría de momentos ni resistencia de compensación de mandos.}

\par\smallskip\noindent{\footnotesize\textit{Datos y solución originales. Las sustituciones registradas se reproducen con SymPy y Studium; eso no establece validación experimental ni revisión experta.}}
